% Adapted from Jake Gutierrez's MIT-licensed resume (see LICENSE).
% Single-column A4 variant. XeLaTeX/Tectonic embeds Unicode fonts.
\documentclass[a4paper,11pt]{article}
\usepackage[margin=1.5cm,top=1.1cm,bottom=1.1cm]{geometry}
\usepackage{fontspec}
\setmainfont{texgyretermes}[Extension=.otf,UprightFont=*-regular,BoldFont=*-bold,ItalicFont=*-italic,BoldItalicFont=*-bolditalic,RawFeature=-kern]
\usepackage{titlesec}
\usepackage{enumitem}
\usepackage[hidelinks]{hyperref}
\pagestyle{empty}
\setlength{\parindent}{0pt}
\setlength{\parskip}{3pt}
\raggedright
\titleformat{\section}{\large\scshape}{}{0em}{}[\titlerule]
\titlespacing*{\section}{0pt}{9pt}{5pt}
\setlist[itemize]{leftmargin=13pt,itemsep=2pt,topsep=3pt,parsep=0pt}
\newcommand{\role}[3]{\textbf{#1}\\#2\\\textit{#3}\par}
\newcommand{\cvcontact}{Madrid, Spain}
\InputIfFileExists{contact.tex}{}{}
\newcommand{\cvheader}[1]{\begin{center}{\Huge\scshape Álvaro Garrido}\\[3pt]#1\\[3pt]\small\cvcontact\end{center}}

\hypersetup{pdftitle={Alvaro Garrido - Engineering CV},pdfauthor={Alvaro Garrido}}
\begin{document}
\cvheader{Aeronautical Engineer | Aeroelasticity \& Loads | Technical Project Management}
\small
\section{Professional Summary}
Aeronautical engineer combining aeroelastic simulation, engineering methodology development and technical project management. Leads a 10-member Aeroelasticity \& Controls team and supervises R\&D objectives. Develops Python/MATLAB tools for load calculations and design optimization, with experience coordinating customers, certification and UAV development.
\section{Experience}
\role{Technical Area Responsible}{Solute | Madrid, Spain}{January 2025 - Present}
\begin{itemize}
\item Manage a 10-member Aeroelasticity \& Controls team, organizing workloads and supervising R\&D objectives to align technical activities.
\item Coordinate customer communications and support commercial requirements, connecting engineering work with customer needs.
\item Prepare technical and commercial proposals to support business development and client portfolio diversification.
\end{itemize}
\role{Senior Project Engineer \& Project Manager}{Solute | Madrid, Spain}{June 2023 - January 2025}
\begin{itemize}
\item Led development of a rotating-wing UAV for wind turbine inspection, covering technical requirements, system architecture, certification processes and testing.
\item Provided technical and functional coordination for wind turbine projects and engineering teams, including life extension and structural integrity.
\item Managed load certification processes against IEC/DNV standards for international projects.
\end{itemize}
\role{Aeroelastic Loads Engineer}{Solute | Madrid, Spain}{December 2021 - June 2023}
\begin{itemize}
\item Developed Python and MATLAB scripts to reduce load calculation cycle times and support wind turbine design optimization.
\item Developed engineering methodologies for Dynamic Wake Meandering and floating wind turbine loads.
\item Performed aeroelastic simulations and dynamic load analyses for GE Renewables offshore turbines, covering bottom-fixed and floating configurations.
\item Conducted site-specific assessments and control tuning for wind turbine projects.
\end{itemize}
\section{Technical \& Management Skills}
\textbf{Methods \& tools:} Python, MATLAB, Git; engineering methodology development; simulation workflow optimization.\\
\textbf{Engineering:} Aeroelasticity, structural dynamics, Bladed, HAWC2; CFD: Fluent; structural tools: NX; CAD: CATIA V5.\\
\textbf{Coordination:} Team leadership, R\&D supervision, customer coordination, technical requirements, system architecture, IEC/DNV load certification.\\
\textbf{Project tools:} MS Project, Jira, Smartsheet.
\section{Education}
\textbf{Master's Degree in Aeronautical Engineering} | 2019 - 2021\\
Universidad Politécnica de Madrid and University of Liège\par
\textbf{Bachelor's Degree in Aerospace Engineering} | 2015 - 2019\\
Universidad Politécnica de Madrid | Average: 7.97/10 | Top 5\% distinction\\
Final project: Carbon fiber recycling using 3D printing - Honors, 9.8/10.
\section{Languages}
Spanish: Native | English: Advanced, professional working proficiency
\end{document}

